% 中文与英文字体微调
\usepackage{fontspec}
\usepackage{xeCJK}

% 标题格式
\usepackage{titlesec}

\usepackage{graphicx}
\usepackage{float}
\usepackage{subcaption}

% 一级标题居中
\titleformat{\section}
  {\centering\Large\bfseries}
  {\thesection}
  {1em}
  {}

% 二级及以下标题靠左
\titleformat{\subsection}
  {\raggedright\large\bfseries}
  {\thesubsection}
  {1em}
  {}

\titleformat{\subsubsection}
  {\raggedright\normalsize\bfseries}
  {\thesubsubsection}
  {1em}
  {}

\setkeys{Gin}{width=0.5\linewidth,keepaspectratio}

\newcommand{\reportimagepair}[5]{%
\begin{figure}[H]
\centering
\begin{subfigure}[t]{0.49\linewidth}
\centering
\includegraphics[width=\linewidth,keepaspectratio]{#1}
\caption{#2}
\end{subfigure}
\hfill
\begin{subfigure}[t]{0.49\linewidth}
\centering
\includegraphics[width=\linewidth,keepaspectratio]{#3}
\caption{#4}
\end{subfigure}
\caption{#5}
\end{figure}
}

\newcommand{\reportsingleimage}[2]{%
\begin{figure}[H]
\centering
\includegraphics[width=0.75\linewidth,keepaspectratio]{#1}
\caption{#2}
\end{figure}
}

% 段落首行缩进
% \setlength{\parindent}{2em}

% 段间距：按中文论文习惯通常不额外拉开
\setlength{\parskip}{0pt}

% 让标题后的第一段也缩进
% \usepackage{indentfirst}

% 代码块长行自动换行
\usepackage{fvextra}
\fvset{
  breaklines=true,
  breakanywhere=true,
  breaksymbolleft={}
}

% 更舒服的行距
\usepackage{setspace}
\setstretch{1.2}
